% ============================================================
\section{Lema C — Asociatividad}
% ============================================================

\begin{frame}{Lema C — Enunciado y el problema}
  \begin{block}{Lema C}
    La multiplicación $\circ$ es asociativa: para cualesquiera
    $\clase{\alpha}, \clase{\beta}, \clase{\gamma} \in \pione(X,\xo)$,
    \[
      \bigl(\clase{\alpha} \circ \clase{\beta}\bigr) \circ \clase{\gamma}
      \;=\;
      \clase{\alpha} \circ \bigl(\clase{\beta} \circ \clase{\gamma}\bigr)
    \]
  \end{block}

  \bigskip
  \textbf{El problema central:} ambos lados recorren los \emph{mismos puntos}
  pero con \emph{distinta repartición de tiempo}:

  \bigskip
  \begin{center}
  \begin{tabular}{lccc}
    \toprule
    Camino & $\alpha$ usa & $\beta$ usa & $\gamma$ usa \\
    \midrule
    $(\alpha*\beta)*\gamma$ &
      $\bigl[0,\tfrac{1}{4}\bigr]$ &
      $\bigl[\tfrac{1}{4},\tfrac{1}{2}\bigr]$ &
      $\bigl[\tfrac{1}{2},1\bigr]$ \\[4pt]
    $\alpha*(\beta*\gamma)$ &
      $\bigl[0,\tfrac{1}{2}\bigr]$ &
      $\bigl[\tfrac{1}{2},\tfrac{3}{4}\bigr]$ &
      $\bigl[\tfrac{3}{4},1\bigr]$ \\
    \bottomrule
  \end{tabular}
  \end{center}

  \bigskip
  La homotopía debe \textbf{desplazar los puntos de corte} suavemente de
  $\bigl\{\tfrac{1}{4},\tfrac{1}{2}\bigr\}$ a $\bigl\{\tfrac{1}{2},\tfrac{3}{4}\bigr\}$
  al ir de $s=0$ a $s=1$.
\end{frame}

\begin{frame}{Lema C — Fórmulas explícitas de los caminos}
  Calculando directamente:

  \[
    \bigl((\alpha*\beta)*\gamma\bigr)(t) =
    \begin{cases}
      \alpha(4t)   & 0 \le t \le \tfrac{1}{4} \\[3pt]
      \beta(4t-1)  & \tfrac{1}{4} \le t \le \tfrac{1}{2} \\[3pt]
      \gamma(2t-1) & \tfrac{1}{2} \le t \le 1
    \end{cases}
  \]

  \bigskip

  \[
    \bigl(\alpha*(\beta*\gamma)\bigr)(t) =
    \begin{cases}
      \alpha(2t)   & 0 \le t \le \tfrac{1}{2} \\[3pt]
      \beta(4t-2)  & \tfrac{1}{2} \le t \le \tfrac{3}{4} \\[3pt]
      \gamma(4t-3) & \tfrac{3}{4} \le t \le 1
    \end{cases}
  \]

  \bigskip
  Los puntos de corte son distintos, pero basta con construir una homotopía
  que los \emph{deslice} continuamente.
\end{frame}

\begin{frame}{Lema C — Construcción de la homotopía $K$}
  Los puntos de corte en el nivel $s$ son $\dfrac{s+1}{4}$ y $\dfrac{s+2}{4}$:

  \[
    \boxed{K(t,s) =
    \begin{cases}
      \alpha\!\left(\dfrac{4t}{s+1}\right)
        & 0 \le t \le \dfrac{s+1}{4} \\[10pt]
      \beta\bigl(4t - 1 - s\bigr)
        & \dfrac{s+1}{4} \le t \le \dfrac{s+2}{4} \\[10pt]
      \gamma\!\left(\dfrac{4t - s - 2}{2 - s}\right)
        & \dfrac{s+2}{4} \le t \le 1
    \end{cases}}
  \]

  \begin{columns}[T]
    \begin{column}{0.48\textwidth}
      \textbf{En $s=0$:} cortes en $\tfrac{1}{4}$ y $\tfrac{1}{2}$
      \[K(\,\cdot\,,0) = (\alpha*\beta)*\gamma \checkmark\]
    \end{column}
    \begin{column}{0.48\textwidth}
      \textbf{En $s=1$:} cortes en $\tfrac{1}{2}$ y $\tfrac{3}{4}$
      \[K(\,\cdot\,,1) = \alpha*(\beta*\gamma) \checkmark\]
    \end{column}
  \end{columns}
\end{frame}

\begin{frame}{Lema C — Verificación de la homotopía}
  \textbf{Punto base fijo:}
  \begin{align*}
    K(0,s) &= \alpha\!\left(\tfrac{0}{s+1}\right) = \alpha(0) = \xo \checkmark\\
    K(1,s) &= \gamma\!\left(\tfrac{4-s-2}{2-s}\right) = \gamma\!\left(\tfrac{2-s}{2-s}\right) = \gamma(1) = \xo \checkmark
  \end{align*}

  \smallskip
  \textbf{Continuidad en las fronteras (Lema de Continuidad):}

  \begin{columns}[T]
    \begin{column}{0.50\textwidth}
      \small $t = \tfrac{s+1}{4}$ \;(frontera $\alpha$/$\beta$):
      \begin{align*}
        \text{p.1: }&\alpha\!\left(\tfrac{s+1}{s+1}\right) = \alpha(1) = \xo \\
        \text{p.2: }&\beta(s+1-1-s) = \beta(0) = \xo \;\checkmark
      \end{align*}
    \end{column}
    \begin{column}{0.50\textwidth}
      \small $t = \tfrac{s+2}{4}$ \;(frontera $\beta$/$\gamma$):
      \begin{align*}
        \text{p.2: }&\beta(s+2-1-s) = \beta(1) = \xo \\
        \text{p.3: }&\gamma\!\left(\tfrac{s+2-s-2}{2-s}\right) = \gamma(0) = \xo \;\checkmark
      \end{align*}
    \end{column}
  \end{columns}

  \smallskip
  Las tres piezas son continuas y coinciden en las fronteras $\Rightarrow$ $K$ es continua.
\end{frame}

\begin{frame}{Lema C — Visualización: los puntos de corte se deslizan}
  \begin{center}
  \begin{tikzpicture}[scale=3.0]
    \draw[thick] (0,0) rectangle (1,1);

    % Regiones coloreadas para s=0 (base)
    \fill[moradoclaro!40] (0,0) rectangle (0.25, 0.01);
    \fill[verdeTeal!30]   (0.25,0) rectangle (0.5, 0.01);
    \fill[coralOscuro!20] (0.5,0) rectangle (1, 0.01);

    % Líneas de corte para distintos valores de s
    % s=0: cortes en 1/4 y 1/2
    \draw[morado, thick]     (0.25, 0) -- (0.25, 0.01)
      node[below=2pt]{\tiny $\frac{1}{4}$};
    \draw[verdeTeal, thick]  (0.5,  0) -- (0.5,  0.01)
      node[below=2pt]{\tiny $\frac{1}{2}$};

    % Líneas diagonales de los puntos de corte
    % Corte izq: t=(s+1)/4  -> de (1/4,0) a (1/2,1)
    \draw[morado, very thick, dashed]
      (0.25, 0) -- (0.5, 1)
      node[right]{\small $\dfrac{s+1}{4}$};

    % Corte der: t=(s+2)/4  -> de (1/2,0) a (3/4,1)
    \draw[verdeTeal, very thick, dashed]
      (0.5, 0) -- (0.75, 1)
      node[right]{\small $\dfrac{s+2}{4}$};

    % Regiones
    \node[morado]     at (0.13, 0.5) {\small $\alpha$};
    \node[verdeTeal]  at (0.48, 0.5) {\small $\beta$};
    \node[coralOscuro] at (0.83, 0.5) {\small $\gamma$};

    % Etiquetas ejes
    \node[below] at (0,-0.07) {\small $0$};
    \node[below] at (1,-0.07) {\small $1$};
    \node[left]  at (-0.04,0) {\small $s=0$};
    \node[left]  at (-0.04,1) {\small $s=1$};

    % Cortes en s=1
    \draw[morado!50]    (0.5, 1) -- (0.5,1.01)
      node[above]{\tiny $\frac{1}{2}$};
    \draw[verdeTeal!50] (0.75,1) -- (0.75,1.01)
      node[above]{\tiny $\frac{3}{4}$};

    % Flecha s
    \draw[-Stealth, thick] (-0.08, 0) -- (-0.08, 1)
      node[midway, left]{\small $s$};
  \end{tikzpicture}
  \end{center}
\end{frame}

\begin{frame}{Lema C — Conclusión}
  \begin{block}{Resultado}
    $K$ es una homotopía módulo $\xo$ entre $(\alpha*\beta)*\gamma$ y
    $\alpha*(\beta*\gamma)$, luego:
    \[
      \bigl[\,(\alpha*\beta)*\gamma\,\bigr]
      \;=\;
      \bigl[\,\alpha*(\beta*\gamma)\,\bigr]
    \]
    y por tanto la operación $\circ$ es \textbf{asociativa}. \hfill $\square$
  \end{block}

  \bigskip
  \textbf{Observación técnica importante:}
  Nótese que $2-s \neq 0$ para $s \in \I$, por lo que $\gamma\!\left(\frac{4t-s-2}{2-s}\right)$
  está bien definida en toda la tercera pieza.  De manera análoga,
  $s+1 \neq 0$ para todo $s \in \I$.
\end{frame}
